\section{Domains on K\"ahler cones}\label{sec:cones}

\subsection{The domain condition}
A singular Riemannian space will mean a locally compact metric space
$X$ with dense smooth regular locus, whose Riemannian length metric
is the restriction of its distance. Forms are complex valued and
are defined on $X_{\reg}$. Their support is taken in $X$, so compact
support may meet the singular set. The maximal operators $d$ and
$\delta$ are defined by distributional differentiation on
$X_{\reg}$; the minimal operator $d_{\min}$ is the graph closure
of $d$ on smooth forms compactly supported in $X_{\reg}$.

We say that the \emph{$L^2$ Stokes property} holds when
$d_{\min}=d_{\max}$. Then
$\delta_{\max}=d_{\max}^*$, and Hilbert complex theory
\cite{brl92} makes
\[
 D=d+\delta,\qquad V(X)=\Dom d_{\max}\cap\Dom\delta_{\max}
\]
self-adjoint, with
$\|Du\|^2=\|du\|^2+\|\delta u\|^2$.
The domain $V(X)$ is called the Gaffney domain. It need not be the
minimal domain of $D$. Bounded Lipschitz functions preserve all the
domains just defined.

Suppose $g$ is K\"ahler on the regular locus. Write $L$ for wedge
product by its K\"ahler form, $\Lambda=L^*$, $W$ for the Weil
operator, and $d^c=W^{-1}dW$, $\delta^c=W^{-1}\delta W$,
$D^c=W^{-1}DW$. The identities on smooth forms are
\begin{equation}\label{eq:kahler-id}
 [L,\delta]=d^c,\quad [\Lambda,d]=-\delta^c,\quad
 [L,d]=[\Lambda,\delta]=0,\quad D^2=(D^c)^2,
\end{equation}
and each of $d,\delta$ anticommutes with each of $d^c,\delta^c$.

\begin{definition}\label{def:WI}
The Gaffney domain is \emph{Weil invariant} on $\Omega\subset X$
if $Wu\in V(X)$ for every $u\in V(X)$ compactly supported in
$\Omega$.
\end{definition}

The condition is local by Lipschitz partitions of unity. It is
equivalent to $d^cu,\delta^cu\in L^2$ for those forms. Since $V$
and $W$ preserve degree, it can be checked on pure-degree forms,
where the two terms of $D^cu$ have different degrees. Assuming
also the Stokes property, on a fixed compact support the closed
graph theorem gives
\begin{equation}\label{eq:WI-estimate}
 \|D^cu\|\le C(\|Du\|+\|u\|).
\end{equation}
Indeed the distributional operator $D^c$ is closed from the
corresponding graph-norm space to $L^2$.

\begin{proposition}\label{prop:global}
If $X$ is compact and satisfies the Stokes property and Weil
invariance, then $d^c\eta=\delta^c\eta=0$ for every
$\eta\in\ker d\cap\ker\delta$. The operators $W,L,\Lambda$
preserve this harmonic space.
\end{proposition}
\begin{proof}
Weil invariance makes $d^c\eta$ and $\delta^c\eta$ square
integrable. They are closed and coclosed by the anticommutation
identities, hence belong to $V$. Applying Weil invariance again
and integrating by parts gives
\[
 \|d^c\eta\|^2+\|\delta^c\eta\|^2
 =\langle\eta,(D^c)^2\eta\rangle
 =\langle\eta,D^2\eta\rangle=0.
\]
All terms in these integrations lie in the required maximal domains;
the Stokes property supplies the adjoint relation. Equation
\eqref{eq:kahler-id} now proves the assertions for $L$ and
$\Lambda$, and the definition proves the assertion for $W$.
\end{proof}

On a hyperk\"ahler regular locus the pointwise operators satisfy
\begin{gather}\label{eq:verbitsky}
 [L_a,L_b]=[\Lambda_a,\Lambda_b]=0,\quad [L_a,\Lambda_a]=h,
 \quad K_{ab}=-K_{ba},\notag\\
 [K_{ab},L_c]=2(\delta_{bc}L_a-\delta_{ac}L_b),\qquad
 [K_{ab},\Lambda_c]=2(\delta_{bc}\Lambda_a-\delta_{ac}\Lambda_b).
\end{gather}
These are Verbitsky's $\mathfrak{so}(4,1)$ relations
\cite{verbitsky90,ll97}. Consequently, Stokes and Weil invariance
for all three structures give this action on harmonic forms.

\subsection{Separation of variables}
Let $Z$ be compact of odd real dimension $f$, and consider the
complete cone at infinity with vertex added,
$C(Z)$, with metric $dr^2+r^2g_Z$. Set $m=f+1$.
We require the Stokes property on $Z$ and on $C(Z)$, and compact
resolvent of $D_Z$ on $V(Z)$. Such a link will be called
\emph{admissible}. A smooth closed odd-dimensional link is
admissible. More generally, the Stokes property on the cone follows
from the first two requirements and oddness of $f$ by Cheeger's
cone theorem; see \cite[\S2]{gl02} for the precise Stokes criterion.

Compact resolvent gives an orthonormal eigenbasis on $Z$, consisting
of harmonic forms and pairs
\[
 \delta_Z\alpha=0,\quad \Delta_Z\alpha=\lambda\alpha,
 \quad \beta=\lambda^{-1/2}d_Z\alpha,
 \quad \delta_Z\beta=\sqrt\lambda\,\alpha\quad(\lambda>0).
\]
The nonzero spectrum is discrete with finite multiplicities.
For degree $j$, write a cone form in unitary radial coordinates as
\[
 u=r^{j-f/2}\varphi(r)+dr\wedge r^{j-1-f/2}\psi(r)
 \equiv(\varphi;\psi).
\]
Its square norm is $\int(\|\varphi\|_Z^2+\|\psi\|_Z^2)\,dr$.
Direct differentiation gives
\begin{align}
 d(\varphi;\psi)&=\big(r^{-1}d_Z\varphi;
 (\partial_r+(j-f/2)/r)\varphi-r^{-1}d_Z\psi\big),\label{eq:d-cone}\\
 \delta(\varphi;\psi)&=\big(r^{-1}\delta_Z\varphi+
 (-\partial_r+(j-1-f/2)/r)\psi;-r^{-1}\delta_Z\psi\big).
 \label{eq:delta-cone}
\end{align}
A homogeneous form has coefficient pair $r^s(a;b)$; it is $L^2$
at zero precisely when $s>-1/2$. The operators $d,\delta$ lower
$s$ by one, whereas $W,L,\Lambda$ and the Hodge star preserve it.

\begin{lemma}\label{lem:indicial}
The homogeneous closed and coclosed forms of degree $j$ with
$L^2$ link coefficients are sums of the following modes:
\begin{enumerate}
\item $r^{s_\pm}(b_\pm\beta;\alpha)$, where
$\deg\alpha=j-1$,
\[
 s_\pm=-\tfrac12\pm\sqrt{\lambda+(j-m/2)^2},\qquad
 b_\pm=\frac{\sqrt\lambda}{s_\pm+j-f/2};
\]
\item $(r^{f/2-j}\eta;0)$ for $\eta\in\cH^j(Z)$;
\item $(0;r^{j-1-f/2}\eta)$ for $\eta\in\cH^{j-1}(Z)$.
\end{enumerate}
In degree $j\ne m/2$ there is no exponent in
$(-1/2-|j-m/2|,-1/2+|j-m/2|)$.
\end{lemma}
\begin{proof}
Equations \eqref{eq:d-cone}--\eqref{eq:delta-cone} require the
tangential component to be closed and the radial component to be
coclosed on $Z$. Their Hodge decompositions therefore leave only
the displayed pairs and the harmonic modes. On a pair, with
$\gamma=j-f/2$, the coefficient equations are
\[
 b'+\gamma b/r-\sqrt\lambda a/r=0,\qquad
 \sqrt\lambda b/r-a'+(\gamma-1)a/r=0.
\]
The indicial equation is $(s+\gamma)(s-\gamma+1)=\lambda$.
The harmonic modes solve the two remaining scalar equations.
Every resulting exponent satisfies $|s+1/2|\ge|j-m/2|$.
\end{proof}

\subsection{The critical modes}
Suppose now the cone is K\"ahler for a dilation-invariant complex
structure and that its Gaffney domain is Weil invariant away from
the vertex. The exponents which can obstruct a cutoff at zero lie
in $-1/2<s\le1/2$.

\begin{lemma}\label{prop:slow}
Every homogeneous closed and coclosed form with exponent in
$(-1/2,1/2]$ and $L^2$ link coefficients is annihilated by both
$d^c$ and $\delta^c$.
\end{lemma}
\begin{proof}
Multiply the form $u$ by a cutoff on an annulus. Weil invariance
there shows that $d^cu$ and $\delta^cu$ have $L^2$ link
coefficients; their exponent is $s-1$.
By Lemma \ref{lem:indicial}, either $u$ has degree $m/2$, or
$s=1/2$ and it is a harmonic-link mode of degree $m/2-1$ or
$m/2+1$.

In the middle degree, $d^cu=-\delta Lu$ and
$\delta^cu=d\Lambda u$ are closed and coclosed. They have
degrees $m/2+1$ and $m/2-1$ and exponent in $(-3/2,-1/2]$,
inside the excluded intervals of Lemma \ref{lem:indicial}.
They vanish.

For degree $m/2-1$, the same argument in degree $m/2-2$ first
gives $\delta^cu=0$. Set $\zeta=d^cu$. It is closed and
coclosed and has degree $m/2$ and exponent $-1/2$. Also
$d^c\zeta=\delta^c\zeta=0$, the latter by
$(D^c)^2u=D^2u=0$. Hence
$W\zeta=r^{-1/2}(\eta_1;\eta_2)$ with harmonic $\eta_i$ on $Z$.
Writing $Wu=r^{1/2}(a;b)$, the equations
$d(Wu)=W\zeta$ and \eqref{eq:d-cone} give
$\eta_1=d_Za$, $\eta_2=-d_Zb$. A harmonic form in the range
of $d_{Z,\max}$ is zero by Stokes on $Z$. Thus $\zeta=0$.
Hodge star reduces degree $m/2+1$ to this case.
\end{proof}

\begin{theorem}[The cone theorem]\label{thm:vertex}
Let $C(Z)$ be a K\"ahler cone with admissible link. Weil invariance
on the complement of the vertex implies Weil invariance at the
vertex as well.
\end{theorem}

For smooth links this is the cone result of
\cite[Theorem~2.1]{cheeger83}; see also
\cite[Theorems~3.7 and~5.9]{brl93}. The singular-link case uses
the preceding lemma and the same scalar radial estimates.
We include those estimates because no spectral gap at the critical
exponents is assumed.

Decompose a pure-degree form into tangential-coexact/radial-exact
terms, tangential-exact/radial-coexact terms, and harmonic-link
terms. Testing weak equations against a link eigenform times a
radial test function is legitimate by Stokes on the cone. It gives
the same equations as \eqref{eq:d-cone}--\eqref{eq:delta-cone}.
Truncations retain both slots of every coupled exact--coexact
block: either slot alone can have a nonintegrable derivative
which cancels only in the pair. The projections onto these
complete blocks preserve $V$ and are orthogonal both for the
$L^2$ norm and the sum of the $d$ and $\delta$ energies.

For the first type, the equations contain the separate terms
$\sqrt\lambda a/r$ and $\sqrt\lambda b/r$, so
\begin{equation}\label{eq:coexact-hardy}
 \|u_{\mathrm{ce}}/r\|\le\lambda_{\min}^{-1/2}\|Du\|.
\end{equation}
For the second type the radial matrix is, up to a unitary sign,
\[
 \partial_r+r^{-1}
 \begin{pmatrix}\gamma&-\sqrt\lambda\\-\sqrt\lambda&1-\gamma\end{pmatrix},
 \qquad \gamma=j-f/2.
\]
Its eigenvalues are
$\theta=1/2\pm\sqrt{\lambda+(j-m/2)^2}$.
Harmonic-link terms give $\theta=\gamma$ or $1-\gamma$.
Thus each coefficient satisfies
\begin{equation}\label{eq:scalar}
 a'+\theta a/r=F.
\end{equation}
The squares of the $L^2$ norms of the right sides sum to the
corresponding part of $\|Du\|^2$.

\begin{lemma}\label{lem:hardy}
Let $a,F\in L^2(0,R)$ satisfy \eqref{eq:scalar}, with $a=0$
near $R$. Then:
\begin{enumerate}
\item for $\theta\ge1/2$ or $\theta<-1/2$,
$\|a/r\|\le|\theta+1/2|^{-1}\|F\|$;
\item for $-1/2<\theta<1/2$,
$a=c\chi r^{-\theta}+a_0$, where $\chi=1$ near zero is a fixed
cutoff and $|c|+\|a_0/r\|\le C_\theta(\|F\|+\|a\|)$;
\item for $\theta=-1/2$, $a=br^{1/2}$ and
$\|b'r^{1/2}\|=\|F\|$.
\end{enumerate}
\end{lemma}
\begin{proof}
On $[\epsilon,R]$, integration by parts gives
\[
 \|a'+\theta a/r\|^2
 =\|a'-a/(2r)\|^2+(\theta+1/2)^2\|a/r\|^2
 +(\theta+1/2)[|a|^2/r]_\epsilon^R.
\]
For $\theta\ge1/2$, integrability forces
$a(r)=r^{-\theta}\int_0^rt^\theta F(t)\,dt$, and the boundary
term tends to zero. For $\theta<-1/2$ it is nonnegative.
For the intermediate range the same integral formula includes
the homogeneous term $cr^{-\theta}$; subtracting $c\chi
r^{-\theta}$ leaves a vanishing boundary term and gives (2).
The constant $c$ is controlled on any fixed interior interval.
The endpoint identity in (3) follows by differentiation.
\end{proof}

\begin{proof}[Proof of Theorem \ref{thm:vertex}]
Estimate \eqref{eq:WI-estimate} on a fixed annulus, dilation,
and a dyadic partition of unity give
\begin{equation}\label{eq:annulus}
 \|D^cw\|\le C(\|Dw\|+\|w/r\|)
\end{equation}
for forms supported away from the vertex. In a fixed degree there
are only finitely many $\theta\in[-1/2,1/2)$, and the other
values have $|\theta+1/2|$ uniformly bounded away from zero.
Equations \eqref{eq:coexact-hardy} and Lemma \ref{lem:hardy}
therefore decompose a compactly supported $u\in V$ as
\[
 u=w+\sum_{-1/2<\theta<1/2}c_\theta\chi u_\theta
       +\sum_{\theta=-1/2}b_\theta u_\theta,
\]
where $u_\theta$ is a closed and coclosed mode of exponent
$-\theta$, $w\in V$, and
$\|Dw\|+\|w/r\|\le C(\|Du\|+\|u\|)$.
All sums are finite. Cut off $w$ at $r=\epsilon$ with derivative
bounded by $C/r$. Estimate \eqref{eq:annulus} is uniform in
$\epsilon$, so distributional convergence gives $D^cw\in L^2$.
Lemma \ref{prop:slow} kills $D^cu_\theta$ in the finite sums.
Their remaining derivative terms are square integrable by the
cutoff estimate or by
$\|b_\theta'r^{1/2}\|<\infty$. Thus $D^cu\in L^2$.
Pure degree separates its two components, proving $Wu\in V$.
\end{proof}

\begin{lemma}\label{lem:compact-cone}
The set of forms supported in $\{r<R\}$ with
$\|u\|+\|Du\|\le1$ is precompact in $L^2(C(Z))$.
\end{lemma}
\begin{proof}
Modes with link eigenvalue $\lambda>\lambda_0\ge4$ have total
$L^2$ norm at most $CR/\sqrt{\lambda_0}$, by
\eqref{eq:coexact-hardy} and Lemma \ref{lem:hardy}(1).
There are finitely many remaining modes. Their coefficients have
bounded $W^{1,2}$ norm on each $[\epsilon,R]$ and uniformly small
mass near zero by Lemma \ref{lem:hardy}. At its endpoint,
$|b(r)|\le\|b'\sqrt r\|\sqrt{\log(R/r)}$ since $b$ vanishes
near $R$. Rellich compactness on the interior intervals, followed
by the uniform tail bounds, proves the assertion.
\end{proof}
